% Einheit 1 von 3 – Irrationale Zahlen und Zahlbereiche (bank/reelle-zahlen/e1.jsonl, zusammenbau v0.1)
%% TODO Zweigzeile Teil 1: was hier gelernt wird (Satz in Schülersprache aus dem Einheitstitel) – Einheit 1
\einheitenkopf[e1]{Einheit 1 von 3 · Irrationale Zahlen und Zahlbereiche}
\zweigzeile{ab Kl. 9 (am Gymnasium ab Kl. 8) · keine P10-Aufgabe}
\setcounter{aufgabe}{9}

%% TODO Titel: Ich-kann-Satz fehlt in der Bank (Platzhalter: Kettenname) – Nr. 10
%% TODO Anweisung: ein Satz über der ersten Teilaufgabe fehlt in der Bank – Nr. 10
\begin{aufgabe}{Zahlbereiche und irrationale Zahlen}
\begin{teile}
\teil $\sqrt{121}$ – geht die Wurzel auf? \\ \kreuz{geht auf} \\ \kreuz{geht nicht auf}
\end{teile}
\begin{teile}
\teil $\frac{5}{8}$ – als Dezimalzahl, abbrechend oder periodisch? Schreibe periodische Zahlen mit Periodenstrich. \\ \kreuz{abbrechend} \\ \kreuz{periodisch} \leerfeld
\end{teile}
\begin{teile}
\teil $\frac{4}{9}$ – als Dezimalzahl, abbrechend oder periodisch? Schreibe periodische Zahlen mit Periodenstrich. \\ \kreuz{abbrechend} \\ \kreuz{periodisch} \leerfeld
\end{teile}
\begin{teile}
\teil $\frac{11}{40}$ – als Dezimalzahl, abbrechend oder periodisch? Schreibe periodische Zahlen mit Periodenstrich. \\ \kreuz{abbrechend} \\ \kreuz{periodisch} \leerfeld
\end{teile}
\begin{teile}
\teil $\frac{5}{6}$ – als Dezimalzahl, abbrechend oder periodisch? Schreibe periodische Zahlen mit Periodenstrich. \\ \kreuz{abbrechend} \\ \kreuz{periodisch} \leerfeld
\end{teile}
\begin{teile}
\teil $\frac{2}{11}$ – als Dezimalzahl, abbrechend oder periodisch? Schreibe periodische Zahlen mit Periodenstrich. \\ \kreuz{abbrechend} \\ \kreuz{periodisch} \leerfeld
\end{teile}
\swfrage{Was bleibt gleich, was ändert sich?}
\begin{teile}
\teil $\sqrt{169}$ – rational oder irrational? \\ \kreuz{rational} \\ \kreuz{irrational}
\end{teile}
%% TODO Darstellung für schwach fehlt in der Bank (reelle-zahlen-e1-k1-s3-v1; Abschnitt „Für schwache Schüler“)
\swz{$\sqrt{5}$ mit dem Taschenrechner – auf drei Stellen gerundet, mit $=$ oder $\approx$?}{}
\swa{$-8$ – in welche Zahlbereiche? Kreuze in der Tabelle an.}{\sachtabelle{ccccc}{Zahl & ℕ & ℤ & ℚ & ℝ}{$-8$ & & & & }}
\swa{$\sqrt{6}$ – markiere die Zahl auf dem Zahlenstrahl. Zwischen welchen Zehnteln liegt sie? \leerfeld $< \sqrt{6} $<$$ \leerfeld}{\zahlenstrahl[xmin=2,xmax=3,xstep=0.1,karo=1]{}}
%% TODO Darstellung für schwach fehlt in der Bank (reelle-zahlen-e1-k1-s6-v1; Abschnitt „Für schwache Schüler“)
\swz{$\sqrt{8}$; $2{,}8$; $2\frac{5}{6}$ – aufsteigend geordnet?}{}
%% TODO Darstellung für schwach fehlt in der Bank (reelle-zahlen-e1-k1-s7-v1; Abschnitt „Für schwache Schüler“)
\swz{Ein Quadrat hat den Flächeninhalt $7\,\text{cm}^2$. Seitenlänge exakt und auf zwei Stellen gerundet?}{}
%% TODO Darstellung für schwach fehlt in der Bank (reelle-zahlen-e1-k1-s8-v1; Abschnitt „Für schwache Schüler“)
\swz{$\sqrt{7} \cdot \sqrt{7}$ – exaktes Ergebnis? Rational oder irrational?}{}
\swa{Einschachtelung von $\sqrt{2}$: Fülle die Tabelle aus. Zwischen welchen Zehnteln liegt $\sqrt{2}$? \leerfeld $< \sqrt{2} $<$$ \leerfeld}{\sachtabelle{cc}{$x$ & $x^2$}{$1{,}4$ & \\ $1{,}5$ & }}
\swz[3]{$\sqrt{40}$: Ist die Zahl rational oder irrational? Begründe. Gib einen Näherungswert auf zwei Stellen an. Zu welchen Zahlbereichen gehört sie?}{}
\end{aufgabe}

%% TODO Titel: Ich-kann-Satz fehlt in der Bank (Platzhalter: Kettenname) – Nr. 11
%% TODO Anweisung: ein Satz über der ersten Teilaufgabe fehlt in der Bank – Nr. 11
\begin{aufgabe}{Dezimalzahl als Bruch schreiben (rational)}
%% TODO Darstellung für schwach fehlt in der Bank (reelle-zahlen-e1-k2-s1-v1; Abschnitt „Für schwache Schüler“)
\swz{$0{,}45$ – als gekürzter Bruch?}{}
\end{aufgabe}

%% TODO Titel: Ich-kann-Satz fehlt in der Bank (Platzhalter: Kettenname) – Nr. 12
%% TODO Anweisung: ein Satz über der ersten Teilaufgabe fehlt in der Bank – Nr. 12
\begin{aufgabe}{Teilmengenkette ℕ ⊂ ℤ ⊂ ℚ ⊂ ℝ erklären (jede natürliche Zahl ist ganz, jede ganze rational, jede rationale reell)}
%% TODO Darstellung für schwach fehlt in der Bank (reelle-zahlen-e1-k3-s1-v1; Abschnitt „Für schwache Schüler“)
\swz[3]{Ist jede ganze Zahl rational? Begründe.}{}
\end{aufgabe}

%% TODO Titel: Ich-kann-Satz fehlt in der Bank (Platzhalter: Kettenname) – Nr. 13
%% TODO Anweisung: ein Satz über der ersten Teilaufgabe fehlt in der Bank – Nr. 13
\begin{aufgabe}{Zahlbereiche und irrationale Zahlen – Fehler finden, Begründen, Darstellungswechsel, Anwendung}
\swz[3]{Lena schreibt $\sqrt{3} = 1{,}732050808$, weil der Taschenrechner das anzeigt. Finde den Fehler und schreibe richtig.}{}
\swfrage{Warum kann $\sqrt{2}$ kein Bruch sein? Erkläre die Idee in Worten.}
\begin{teile}
\teil Am Zahlenstrahl ist der Punkt A markiert. Welche Zahl ist A? \\ \kreuz{$\sqrt{5}$} \\ \kreuz{$\sqrt{8}$} \\ \kreuz{$\sqrt{11}$}
\end{teile}
\swz[3]{Ein quadratischer Teppich soll genau $6\,\text{m}^2$ groß sein. Wie lang muss eine Seite sein? Gib den exakten Wert und einen Wert für den Zuschnitt an.}{}
\end{aufgabe}

\uebersichtskasten{Rational: Jede Zahl, die man als Bruch schreiben kann, ist rational – als Dezimalzahl bricht sie ab oder wird periodisch: 3/8 = 0,375; 1/3 = 0,333… \\ Irrational: Eine Zahl, die weder abbricht noch periodisch wird, ist irrational – sie hat keinen Bruch, nur Näherungswerte: √2 = 1,41421…, π = 3,14159… \\ Wurzel rational oder nicht: Ist der Radikand eine Quadratzahl, geht die Wurzel auf und ist rational, sonst ist sie irrational: √49 = 7, √50 ≈ 7,07. \\ Zahlbereiche: ℕ ⊂ ℤ ⊂ ℚ ⊂ ℝ – die reellen Zahlen sind rationale und irrationale Zahlen zusammen; jeder Punkt der Zahlengeraden gehört zu genau einer reellen Zahl.}
